\documentclass[11pt]{article}
\usepackage[T1]{fontenc}
\usepackage{amsmath,amssymb,amsthm}
\usepackage[margin=1in]{geometry}
\usepackage[hidelinks]{hyperref}
\newtheorem{theorem}{Theorem}
\newtheorem{lemma}[theorem]{Lemma}
\newtheorem{proposition}[theorem]{Proposition}
\newtheorem{corollary}[theorem]{Corollary}
\newcommand{\C}{\mathcal C}
\newcommand{\D}{\mathcal D}
\newcommand{\RR}{\mathbb R}
\newcommand{\aa}{\mathfrak a}
\title{Normalization compatibility for geometric obstacle envelopes\\
Conductor localization, full fibres, and quantitative discontinuity}
\author{Independent scoped ORIGINAL P12.3 deductions; unreviewed draft}
\date{9 October 2026}
\begin{document}
\maketitle
\begin{abstract}
For a finite normalization of an excellent local domain we identify the
raw geometric envelope with the normal-source envelope at every point
having finite conductor value. If the normal-source envelope is continuous
and that locus is dense, a greatest target minorant exists exactly when
the normal envelope is constant on every full restriction fibre.
A discrepancy of size $\Delta$ on one fibre gives a non-upper-semicontinuous
raw supremum, a uniform error at least $\Delta$ for every canonical degree,
and the same gap for every decreasing-completed minorant. We prove the
needed full-support density for smooth closed source germs by actual
translated Gauss seminorms on formal coordinate discs. This yields an
obstacle-by-obstacle criterion for arbitrary conductor geometry whenever
the normalization centres are smooth affine closed germs over an
algebraically closed field of characteristic zero. It does not classify
all rings by conductor injectivity or close the unchanged original problem.
\end{abstract}

\section{Geometric classes and exact hypotheses}
Let $R$ be an excellent local domain of positive dimension, $S$ its finite
normalization, a normal semilocal domain, and $\mathfrak c$ its nonzero
conductor. Thus $\mathfrak c$ is an $S$-ideal contained in $R$.
Let $\aa\subset R$ be proper primary and put
\[
 X=V_{\aa}(R),\qquad Y=V_{\aa S}(S),\qquad \rho:Y\longrightarrow X.
\]
Every marked normalization centre, every prime support, infinite values
of individual elements, and the full compact pointwise topology are used.
The entire qualified G331 gives the continuous surjective compact quotient
$\rho$, a singleton fibre wherever $v(\mathfrak c)<\infty$, and its complete
conductor-quotient description on the remaining fibres.
Write
\[
 Y^\circ=\{w:w(\mathfrak c)<\infty\},\quad
 X^\circ=\{v:v(\mathfrak c)<\infty\}.
\]
The values here use the same $R$-elements of the conductor; source and
target conductor values agree.

The cones $\C,\D$ are G331's independently geometric continuous uniform
cone and actual decreasing-$\C$-net completion. Positivity is actual
negative effective finite real Cartier presentations on directed actual
primary blowups, with nonnegative degree on every contracted proper
integral curve. We use G331's proved comparison with finite negative ideal
profiles, not a substitute definition. Pullback preserves these cones;
full-fibre-constant source members descend positively. Proper primary
profiles have actual relatively ample presentations before any limit;
the unit profile is zero.

For continuous signed $f$ on $X$ let $F=\rho^*f$. The raw suprema over
$\C$ and $\D$ minorants are denoted $p_R(f)$ and $p_S(F)$.
The canonical lemma recalled below proves they coincide within each ring.
They are finite and nonpositive. No envelope regularity for $R$ or $S$
is assumed in the next localization lemma.

\section{The canonical lemma, including every support and degree}
For either ring and its section define
\[
 A_m(f)=\{s:v(s)+mf(v)\ge0\text{ for all }v\},\qquad
 h_m(f)=-v(A_m(f))/m,\qquad m\ge1.
\]
Ultrametricity makes these ideals. A sufficiently large boundary power
lies in each, so they are finitely generated and unit-or-primary.
Each $h_m$ is an actual continuous geometric minorant, and
$A_mA_n\subset A_{m+n}$. Infinite values already satisfy the tests.

If $g\in\C$, $g\le f$, take a rational primary profile
$q=-v(I)/k$ with $\|q-g\|<\epsilon$. Then
$I^N\aa^{\lceil Nk\epsilon\rceil}\subset A_{Nk}(f)$ and
$h_{Nk}\ge q-\lceil Nk\epsilon\rceil/(Nk)$.
Let $N$ tend to infinity and then $\epsilon$ decrease.
If $g=\inf_i g_i\in\D$, $g\le f$, the directed open sets
$\{g_i<f+\epsilon\}$ cover the compact section. One common index from a
finite subcover makes $g_i-\epsilon$ a $\C$ minorant dominating
$g-\epsilon$, including partial bottom. Consequently
\[
 p_C(f)=p_D(f)=\sup_m h_m(f).
\]
Writing $M_1=\|v(A_1(f))\|_\infty$, grading for $n=km+r$, $0\le r<m$,
gives $h_n\ge(km/n)h_m-rM_1/n\ge h_m-mM_1/n$.
Hence every degree tends pointwise to the raw supremum.
If it is continuous, finitely many strict tests $h_m>p-\epsilon$ cover
the section; their common multiple $N$ has $h_N>p-\epsilon$ everywhere.
Grading yields $0\le p-h_n\le\epsilon+NM_1/n$ for late $n$.
The resulting uniform limit belongs to $\C$ and is the greatest minorant.
Conversely a greatest $\D$ minorant is this raw supremum, both lower
semicontinuous as a supremum of continuous functions and upper
semicontinuous as an actual decreasing completion. It is continuous.
Negative gauge shifts also give
\[
 \|p(f)-p(f')\|_\infty\le\|f-f'\|_\infty .
\]
No finite generation or source resolution is used in this lemma.

\section{Conductor-primary competitors and exact localization}
\begin{lemma}\label{local}
On all of $Y$, $\rho^*p_R(f)\le p_S(F)$. On $Y^\circ$ there is equality.
The equality holds without assuming either raw envelope continuous.
\end{lemma}
\begin{proof}
Pullback of every target competitor is a source competitor, giving the
first inequality. Fix a source canonical ideal $J=A_m^S(F)$ and its
profile $q=-w(J)/m\le F$. Choose an integer
$M\ge\max(\|q\|_\infty,\max(0,-\min f),1)$.
For $n\ge1$ set the actual target ideal
\[
 B_n=\mathfrak c J^n+\aa^{mnM}\subset R.
\]
The first summand is an $R$-ideal because $\mathfrak c$ is an $S$-ideal
contained in $R$. It is a finite $R$-module, hence finitely generated.
The second summand makes $B_n$ unit-or-primary. At every source point
the profile is
\[
 -\frac{w(B_n)}{mn}
 =\max\left(q(w)-\frac{w(\mathfrak c)}{mn},-M\right)\le F(w).
\]
Indeed values of products and sums of finitely generated ideals give
this identity, also when $w(\mathfrak c)=\infty$. Computing a finite
$R$-generating set or a finite $S$-generating set gives the same ideal
value: coefficients from either ring have nonnegative values.
The identity uses only actual ring elements; no infinite supported
semivaluation is extended to $\operatorname{Frac}(S)$.
Surjectivity of $\rho$ makes $-v(B_n)/(mn)$ a target minorant of $f$.
Its proper case has an actual negative effective relatively ample
primary presentation, positive on every contracted integral curve.

At $w\in Y^\circ$, these profiles tend to $q(w)$ as $n$ increases,
since $q(w)\ge-M$. Thus $\rho^*p_R(f)(w)\ge q(w)$.
Take the supremum over all $m$ and apply the canonical lemma to the
source. This proves equality there, without giving positivity to an
increasing raw supremum or assuming finite generation.
\end{proof}

\section{A full-fibre compatibility criterion}
\begin{theorem}\label{criterion}
Suppose $Y^\circ$ is dense in $Y$ and $p_S(F)$ is continuous.
A greatest $\C_R$ or $\D_R$ minorant of $f$ exists if and only if
$p_S(F)$ is constant on every full fibre of $\rho$.
When this holds the greatest minorant is continuous, lies in $\C_R$,
and pulls back exactly to $p_S(F)$.
\end{theorem}
\begin{proof}
If the source envelope is fibre-constant, the compact Hausdorff quotient
descends it to a continuous $g$ on $X$. Source continuity and the canonical
lemma put $p_S(F)$ in $\C_S$; G331's full-fibre positive descent puts
$g$ in $\C_R$. It is below $f$, and Lemma \ref{local} bounds every target
competitor by it. It is therefore the greatest minorant in both classes.

Conversely any target greatest minorant equals the raw supremum and is
continuous by the canonical lemma. Its pullback is continuous and,
by Lemma \ref{local}, equals the continuous $p_S(F)$ on the dense set
$Y^\circ$. They agree everywhere, so the source envelope is fibre-constant.
This argument does not replace a positive descent theorem by ordinary
descent of continuous values, or assign positivity to a fibre minimum.
\end{proof}

\begin{proposition}\label{gap}
Under the same hypotheses, suppose $\rho(w_-)=\rho(w_+)=v$ and
$p_S(F)(w_+)-p_S(F)(w_-)=\Delta>0$. Then
\[
 \limsup_{x\to v}p_R(f)(x)-p_R(f)(v)\ge\Delta.
\]
Every $\D_R$ minorant $g\le f$, including every canonical $h_n$, satisfies
$\sup_X(p_R(f)-g)\ge\Delta$ in the extended sense.
There is no greatest minorant, and no finite generation of the canonical
graded algebra. If source universal continuity is known, perturbing $f$
by less than $\Delta/2$ preserves failure, with gap at least
$\Delta-2\|f-f'\|_\infty$.
\end{proposition}
\begin{proof}
By the full-source bound,
$p_R(f)(v)\le p_S(F)(w_-)$. Choose a net $w_i\in Y^\circ$ tending to
$w_+$; its restrictions $v_i$ tend to $v$. Localization and source
continuity give $p_R(f)(v_i)\to p_S(F)(w_+)$.
This proves the non-upper-semicontinuity and its stated size.
Every $\D$ member is upper semicontinuous, so
$\limsup_i g(v_i)\le g(v)\le p_S(F)(w_-)$.
The uniform gap follows even if some values are bottom.
Finite homogeneous generation would make the raw envelope a rational
primary profile: if homogeneous generators $b_j$ have degrees $d_j$,
every degree has normalized value at least $\min_j v(b_j)/d_j$ and their
powers attain the infimum. With $L$ a common multiple,
$p=-v((b_j^{L/d_j}))/L$, a unit-or-primary continuous profile.
Indeed degree-one generators generate $A_1$, which contains a boundary
power. Their powers occur in the displayed ideal, so its radical is
maximal in the proper case; if it contains a unit the profile is zero.
This contradicts the proved discontinuity.
Finally the canonical Lipschitz bound changes each of the two source
values by at most $\|f-f'\|_\infty$. Source universal continuity supplies
the hypotheses for $f'$, giving the perturbation assertion.
\end{proof}

\section{Density on smooth closed normalization centres}
\begin{lemma}\label{density}
Let a marked source local ring be a smooth closed germ over an
algebraically closed field $k$. For every nonzero conductor ideal,
its finite-value locus is dense in the full normalized section,
including points with nonminimal support. Characteristic zero is not
needed for this density assertion.
\end{lemma}
\begin{proof}
Smooth local parameters give $\widehat S=k[[z_1,\ldots,z_d]]$.
For completeness, the formal-coordinate assertion follows by choosing
an affine presentation with an invertible Jacobian minor and recursively
solving each coefficient of the remaining variables; the invertible
linear system has a unique solution in every degree. It gives inverse
maps to the completed parameter ring. No singular Cohen presentation
is used. G331 extends a normalized semivaluation $w$ to the completion
and contracts primary ideals exactly.

Quotient by its support, form the valued quotient fraction field and
complete it to $H$. Let $\alpha_i$ be the image of $z_i$ in $H$.
All $|\alpha_i|<1$, and $k$ embeds with trivial absolute value.
Choose $\max_i|\alpha_i|<\sigma<1$. Every formal series over $k$
converges in the radius-$\sigma$ Tate algebra over $H$, with its original
nonzero coefficients still nonzero. Translation $z_i=\alpha_i+T_i$ is
an automorphism of that algebra, with inverse $T_i=z_i-\alpha_i$;
both substitutions converge since $|\alpha_i|<\sigma$.
For $0<r<\sigma$, restrict the translated series to radius $r$ and
take its Gauss norm. Restriction is injective because the same power-series
coefficients remain, and translation is invertible. Thus every nonzero
formal series has strictly positive norm and finite associated value.
The weighted Gauss norm is multiplicative: a least-weighted term exists,
and its nonzero initial polynomial belongs to the polynomial ring over
the associated graded valued field, with variables assigned weight
$-\log r$. That graded ring is a domain, so initial polynomials have
nonzero product. Convergent truncations then give the series assertion.

These norms define actual seminorms $w_r$ on the source with zero support.
They are nonnegative on the formal ring and positive on its maximal
ideal. As $r$ decreases to zero, the norm of a translated series tends
to the norm of its constant coefficient, its value at $\alpha$.
To see this on every series, its radius-$\sigma$ norm bounds the
positive-degree radius-$r$ terms by a constant times $r/\sigma$.
Hence $w_r(s)\to w(s)$ in $[0,\infty]$, also for elements killed by $w$.
For finitely many generators of the boundary ideal,
$w_r(\aa S)\to w(\aa S)=1$. Divide $w_r$ by this positive value;
the resulting actual normalized points still converge on every source
element and have finite conductor value. All formal-series and pointwise
limits have been proved; no test only on a finite coordinate list or
only on divisorial points suffices here.
Apply this argument independently at each of the finitely many clopen
source centres. Their union proves density on $Y$.
\end{proof}

\begin{corollary}
Suppose the finite normalization centres are smooth affine closed germs
over an algebraically closed characteristic-zero field.
For EVERY continuous signed target obstacle, greatest geometric minorants
exist exactly when the normal-source envelope is constant on every full
restriction fibre. Any discrepancy gives the exact lower bound of
Proposition \ref{gap} for every degree and is stable under small perturbation.
Universal target continuity is equivalent to this compatibility for all
pulled target obstacles.
\end{corollary}
The source regularity is precisely the qualified BFJ smooth affine CLOSED
domain, applied on the finite product of marked centres. The finite product
envelope is the product of their envelopes. Lemma \ref{density} and
Theorem \ref{criterion} give the result. This allows arbitrary same-centre
or different-centre conductor geometry. It does not say that every
noninjective conductor boundary fails universally: the compatibility of
the normal envelope remains the actual criterion.

\section{Source qualifications and original residual}
The sole correctness grant remains ENTIRE qualified G331, ALL FIVE fields
and mandatory model/support/radial/source/historical-audit qualifications.
Exact full normalization fibres, primary comparison, positive full-fibre
descent, completion and gauges have been rejudged. The entire current426n1
canonical, surface and smooth proof was freshly physically reread for this
application. The three complete BFJ setting, simple-modification and
3.15--3.25 immediate units retain the fresh ROOT428 reading and exact
application, with no extra primary reading credit for this paper.
The real parameter $t=m+C$, thinness bounded on a FIXED dual complex,
substantive closed-simplex continuity and infima over ALL retractions
remain necessary in the smooth source proof. Historical principalization,
multiplier subadditivity and Izumi are published background, not a complete
new recursive audit. Primary literal bodies and pixels remain local.

For other normal sources the criterion is conditional on the two stated
hypotheses: source envelope continuity and density of the finite-conductor
locus. No density or higher-dimensional regularity is inferred merely
from normality. Exact G331 completion carries proved cases and primary
contraction; another primary gauge divides profiles, obstacles and gaps
by the same bounded positive fibre constant. In particular a fibre gap
$\Delta$ becomes $\Delta/d(v)$ and nonexistence persists.
If $R$ is normal, conductor is unit and fibres are singletons, so the
criterion reduces to its source hypothesis. Empty sections are direct.

An intrinsic ring classification for arbitrary conductor covers, general
excellent NORMAL dimension at least three, higher-dimensional positive-
or mixed-characteristic source regularity, and the full original P12.3
remain OPEN. Products, energy, full operators and literal radial comparison
are separate targets. This proof was developed independently while
native296297 were active, with no mathematical inspection or hot feedback.
Compilation and acceptance are unverified; candidate is false.
\end{document}
